% ===== §3 Power Structures and Wealth in Intimate Relations =====
\section{Power Structures and Wealth in Intimate Relations}
\label{p15:sec:power}

\noindent
The preceding section mapped the major traditions of wealth theory and identified their shared individualist presupposition. This section turns to the concrete question of how wealth forms organised by individualist and mercantilist logic actually operate within intimate relational fields---how they enter the field, what they do to the relational dynamics, and through what specific mechanisms they generate or destroy the conditions for GRW generation.

The analysis is not a moral indictment of any particular cultural practice. The historical and material conditions under which traditional wealth forms developed, and the functions they served within those conditions, are addressed in \xref{p15:sec:hegel_marx} under the framework of historical materialism. The present section is concerned with the structural analysis: how do specific wealth forms interact with the relational field's co-evolutionary dynamics? The answer, we argue, turns on three mechanisms---settlement, commodification, and indebtedness---that operate, in different combinations and intensities, in both traditional and modern wealth forms in intimate relations.

\subsection{Bride-Price and Gold: The Mercantilist Form in Intimate Relations}
\label{p15:subsec:brideprice}

\noindent
The institution of bride-price---the transfer of material wealth from the groom's family to the bride's family as a condition of marriage---is among the most widely distributed cultural practices in human history, appearing in various forms across sub-Saharan Africa, East and Southeast Asia, the Middle East, and parts of Europe and the Americas \parencite{mbiti1970}. Its contemporary forms range from the explicit payment of livestock or cash to the more implicit expectation of gifts, hospitality expenditure, and material demonstration of the groom's capacity to support a household.

It is important, before the structural analysis, to acknowledge the diversity of meaning and function that bride-price carries across cultural contexts. In many traditions, bride-price is not understood as a purchase price for the woman but as a gift to her family expressing gratitude and establishing alliance; as a material guarantee of the groom's commitment and capacity; as a mechanism for redistributing wealth across family networks; or as a form of social insurance that provides the bride's family with resources in the event of the marriage's dissolution. These functions are real and should not be dismissed in favour of a reductive reading of bride-price as simple commodity exchange.

Nevertheless, the structural analysis reveals that bride-price, whatever its cultural framing, operates through mechanisms that have determinate effects on the intimate relational field. The question is not what bride-price \emph{means} to the parties involved but what it \emph{does} to the relational dynamics.

\emb{What bride-price does} is introduce a pre-existing exchange structure into the intimate relational field---a structure that precedes the relation, is determined by factors external to the relational dynamics (the families' negotiations, the cultural norms for appropriate payment, the assessment of the woman's value by criteria external to the relation itself), and that creates a set of relational obligations that are not generated by the co-evolutionary activity of the intimate field but imposed upon it from outside.

This has three specific structural consequences for the relational field. First, it establishes a \emb{symbolic account} of the relational value that has been transferred, making the relation legible to the symbolic order as a completed transaction. The relational value is, in this sense, pre-settled: it has been assessed, quantified, and exchanged before the intimate relation properly begins. Second, it creates a \emb{power asymmetry} within the relational field: one party (the bride) has been the object of the exchange rather than its subject, and this asymmetry is encoded in the founding symbolic structure of the relation. Third, it generates a \emb{baseline of expectation} about the character of relational exchange---if the founding transaction of the relation was a commodity exchange, then subsequent exchanges within the relation are more likely to be framed in terms of exchange logic rather than co-evolutionary generativity.

None of these consequences is inevitable or irreversible. Intimate relations that begin with bride-price can and do develop rich co-evolutionary dynamics that generate GRW; the founding structure does not determine the trajectory. But the structural tendency is real: bride-price introduces mercantilist logic into the founding moment of the intimate relation, and this founding moment shapes the symbolic structure within which subsequent relational dynamics unfold.

\subsection{Three Mechanisms of Relational Wealth Destruction}
\label{p15:subsec:threemechanisms}

\noindent
The analysis of bride-price reveals three mechanisms that operate, in various combinations, across traditional and modern forms of wealth organisation in intimate relations. We call these settlement, commodification, and indebtedness.

\subsubsection{Settlement: The Closure of Relational Value}
\label{p15:subsubsec:settlement}

\noindent
Settlement (\zh{结算化}) is the mechanism by which relational value is converted into a fixed, symbolic quantity---a price, a score, a balance---that closes the relational account. Settlement transforms relational value from an open, generative process into a closed, enumerable fact.

The paradigmatic form of settlement in intimate relations is the assignment of a monetary value to the relation itself: the bride-price that assesses the woman's relational value as a quantity of gold; the divorce settlement that converts the shared life into a balance of assets to be divided; the prenuptial agreement that pre-settles the distribution of relational value before the relation begins. In each case, the living, dynamic, open-ended process of co-evolutionary value generation is frozen into a snapshot---a moment at which the accounts are balanced and the value is fixed.

Settlement is destructive of GRW for a reason that follows directly from the GRW concept: GRW is constitutively non-settleable. Its value cannot be captured in any fixed quantity because it is not a stock but a process---the ongoing co-evolutionary dynamics of the relational field. To settle GRW is to kill it: the moment the relational value is fixed as a quantity and the account is closed, the co-evolutionary process that generates GRW is interrupted.

This does not mean that all formal agreements in intimate relations are settlement in the destructive sense. The distinction that \pinkref{Paper~XIII} drew between the non-binding vow and the sword of legal enforcement is relevant here \parencite{paperxiii}: an agreement that establishes the conditions for co-evolutionary activity without fixing its products is structurally different from an agreement that converts the products of co-evolutionary activity into a settled account. The former supports GRW generation; the latter destroys it.

\subsubsection{Commodification: The Reduction of the Other to Exchange Value}
\label{p15:subsubsec:commodification}

\noindent
Commodification (\zh{商品化}) is the mechanism by which the relational other---the person with whom one builds a shared life---is reduced to an object with exchange value: a bundle of properties that can be assessed, compared, and traded. Commodification is the intimate relation's form of Marx's commodity fetishism: the irreducibly particular other, with their unique history and their specific mode of co-evolutionary participation, is reduced to a set of general, exchangeable properties.

The most explicit forms of commodification in intimate relations are the assessment of potential partners in terms of their economic value (income, assets, earning potential), their reproductive value (fertility, health, genetic endowment), or their social value (family connections, educational credentials, status). Each of these assessments reduces the other to a bundle of properties that can be compared across potential partners, as commodities are compared across markets.

But commodification operates in subtler forms within ongoing intimate relations. When one partner treats the other's emotional labour as a service to be extracted rather than a co-evolutionary contribution to be received; when the shared life is assessed in terms of its return on investment rather than its co-evolutionary depth; when the other is valued for what they provide rather than for what they are---commodification is operating, even without explicit monetary exchange.

Commodification is destructive of GRW because GRW requires the recognition of the other as a genuine co-subject of the relational field---as someone whose particular, irreducible mode of co-evolutionary participation is constitutive of the field's dynamics. The commodified other cannot be a genuine co-subject; they are a resource. And resources can be extracted but not co-evolved with.

\subsubsection{Indebtedness: The Structural Obligation That Forecloses Generativity}
\label{p15:subsubsec:indebtedness}

\noindent
Indebtedness (\zh{债务化}) is the mechanism by which a structural obligation---a debt that one party owes to another---is introduced into the intimate relational field. Unlike settlement (which closes an account) and commodification (which reduces the other to exchange value), indebtedness keeps the account open---but in a way that is destructive rather than generative.

The paradigmatic form of indebtedness in intimate relations is the bride-price debt: the groom's family has made a significant material transfer to the bride's family, and this creates an implicit or explicit obligation on the bride's part to fulfil the terms of the exchange---to bear children, to maintain the household, to provide the services that the transfer was understood to purchase. The debt is structural: it does not need to be explicitly invoked to shape the relational dynamics; it is built into the founding symbolic structure of the relation.

But indebtedness is not confined to bride-price arrangements. It operates wherever one party's contribution to the intimate relation---material, emotional, or otherwise---is framed as a loan rather than a gift or a co-evolutionary investment: wherever the implicit message is ``I have given you X, and you owe me Y.'' This framing is destructive of GRW because it converts co-evolutionary generativity into exchange: the contributor's investment in the relational field is not received as a contribution to the shared attractor landscape but as a claim on the other's future behaviour.

The deepest form of indebtedness is not material but existential: the sense that one party owes their very existence, happiness, or security to the other, and that this debt can never be fully repaid. This existential indebtedness---which can arise from material dependency, from asymmetric emotional investment, or from cultural structures that assign one party responsibility for the other's wellbeing---is perhaps the most powerful destroyer of genuine co-evolutionary dynamics, because it converts the open, generative space of the intimate relation into a closed system of obligation and repayment.

\begin{claim}[Three mechanisms of GRW destruction]
Settlement, commodification, and indebtedness are three mechanisms by which individualist and mercantilist wealth logic destroys the conditions for GRW generation in intimate relations. Settlement closes the relational account, fixing as a snapshot what is constitutively a process. Commodification reduces the co-subject of the relational field to a bundle of exchangeable properties, foreclosing genuine co-evolutionary participation. Indebtedness converts co-evolutionary generativity into a system of obligation and repayment, replacing the open space of shared becoming with the closed structure of debt.
\end{claim}

\subsection{Social Reproduction and the Invisible Wealth of Care}
\label{p15:subsec:socialreproduction}

\noindent
Federici's analysis of social reproduction \parencite{federici2004} identifies a further dimension of wealth destruction in intimate relations: the systematic non-recognition of care labour as wealth production. We develop this analysis here in terms of the GRW framework, showing how the non-recognition of care labour is not merely an injustice but a form of GRW destruction.

Care labour---the emotional, attentional, and physical work of maintaining the intimate relational field---is GRW's primary production mechanism. The work of joint attention, of timely sharing, of co-evolutionary receptivity, of the non-possessive participation in the relational field's generative activity that \pinkref{Paper~XIV} identified as the practices of relational happiness \parencite{paperxiv}: all of this is care labour in Federici's sense. It generates GRW---the accumulated attractor landscape of the shared life---but it generates it in a form that is invisible to the dominant wealth frameworks.

This invisibility is not innocent. When care labour is not recognised as wealth production, two consequences follow. First, the care labourer is deprived of the social recognition and material support that would allow them to sustain and develop their care practice---creating the material conditions under which care labour is degraded and the GRW it generates is depleted. Second, the value generated by care labour is available to be appropriated by the non-caring party without recognition or reciprocity---a form of GRW extraction that is structurally identical to the extraction of surplus value in Marx's account of capitalist production.

The recognition of care labour as GRW production is therefore both a matter of epistemic justice (recognising what is genuinely there) and a matter of GRW sustainability (ensuring that the conditions for GRW generation are maintained).

\subsection{The Intergenerational Reproduction of Wealth Structures}
\label{p15:subsec:intergenerational}

\noindent
Thomas Piketty's analysis of wealth inequality \parencite{piketty2014} identifies a structural tendency of capital to reproduce and amplify itself across generations: when the rate of return on capital exceeds the rate of economic growth, wealth inequality increases over time, and the advantages of those who begin with wealth compound across generations.

This analysis applies to the intergenerational reproduction of wealth structures in intimate relations. The material wealth that a family brings to the founding of a new intimate relation---the bride-price, the dowry, the parental financial support, the inheritance---shapes the symbolic structure within which the new relation develops. Families with greater material wealth can provide their children with more favourable entry conditions into the intimate relational market; families with less material wealth are constrained in ways that limit their children's relational possibilities.

But the intergenerational reproduction of wealth structures in intimate relations is not only material. The \emb{symbolic capital} that Bourdieu would call the habitus of relational life---the implicit knowledge of how to navigate intimate relations, the emotional and attentional skills developed in the family of origin, the capacity for co-evolutionary participation that is cultivated or suppressed in childhood---is also reproduced intergenerationally. Children who grow up in relational fields rich in GRW---where care labour is valued, where emotional attunement is practised, where co-evolutionary generativity is the norm---develop the relational capabilities that make GRW generation possible in their own intimate relations. Children who grow up in relational fields organised by settlement, commodification, and indebtedness develop different habits and expectations.

This intergenerational dimension of relational wealth has a disturbing consequence: the conditions for GRW generation are themselves unequally distributed. The children of relational fields rich in GRW are better equipped to generate GRW in their own relations; the children of relational fields organised by mercantilist logic are more likely to reproduce that logic. This is a form of relational inequality that is invisible to standard measures of wealth distribution but that is, in the long run, perhaps the most consequential form of inequality in human social life.

\subsection{Contemporary Forms: Consumerism, Human Capital, and Digital Extraction}
\label{p15:subsec:contemporary}

\noindent
The mercantilist logic that operates in traditional wealth forms in intimate relations has not disappeared in contemporary societies; it has taken new forms that are in some respects more pervasive and more insidious than the traditional ones, because they operate through cultural norms and technological systems rather than through explicit contractual structures.

\emb{Consumerism} introduces settlement and commodification into intimate relations through the medium of consumption. The commodification of romantic love---the expectation that love should be expressed through the purchase of goods, that the quality of the relation is measured by the quality of the gifts exchanged, that the anniversary dinner at an expensive restaurant is the proper form of relational celebration---converts the co-evolutionary dynamics of the intimate field into a series of market transactions. The happiness jar, in this cultural environment, is replaced by the receipt: the evidence that the appropriate commodity has been purchased and the relational account has been settled for another year.

\emb{Human capital theory}, associated with the work of Gary Becker, introduces the language of investment, return, and efficiency into the analysis of intimate relations \parencite{maslow1954}. Partners are assessed as human capital assets; the decision to marry is modelled as an investment decision; the allocation of time between market work and domestic labour is optimised for household income maximisation. This framework imports the full apparatus of commodity exchange logic into the most intimate dimensions of human life, converting the co-subject of the relational field into a bundle of productive capacities to be assessed and utilised.

\emb{Digital platforms} have created a new form of relational wealth extraction that is unprecedented in its scale and its invisibility. The data generated by intimate relational activity---the messages, the shared photographs, the location data, the emotional expressions---is harvested by platform companies and converted into advertising revenue. This is a form of GRW extraction: the value generated by co-evolutionary activity in the intimate relational field is captured by an external party without recognition or reciprocity. The platform does not merely observe the intimate relation; it shapes it, through algorithmic design that optimises for engagement rather than for co-evolutionary depth, thereby actively degrading the conditions for GRW generation while extracting the value that the degraded relation still produces.

\begin{claim}[Contemporary forms of GRW extraction]
Contemporary consumerism, human capital theory, and digital platform extraction represent new forms of the mercantilist logic that operates in traditional wealth forms in intimate relations. Each converts the co-evolutionary dynamics of the intimate field into a system of market transactions, reducing the co-subject to a bundle of exchangeable properties and extracting the value generated by relational activity without recognition or reciprocity. The scale and invisibility of digital platform extraction make it perhaps the most significant contemporary threat to GRW generation.
\end{claim}
